%----------------------------------------------------------------------------------------
%    PACKAGES AND THEMES
%----------------------------------------------------------------------------------------

\documentclass[aspectratio=169,xcolor=dvipsnames]{beamer}
\usetheme{SimpleDarkBlue}
% --- Fix algorithm font corruption / ligatures / characters ---
\usepackage[T1]{fontenc}
\usepackage{lmodern}
\usepackage{textcomp}
\usepackage{microtype}
\usepackage{amssymb}

\setbeamertemplate{bibliography item}{\insertbiblabel}


\usepackage{cmap}
\usepackage{hyperref}
\usepackage{graphicx} % Allows including images
\usepackage{booktabs} % Allows the use of \toprule, \midrule and \bottomrule in tables
\usepackage{tikz} 
\usetikzlibrary{shapes.geometric,arrows.meta}
\usepackage{algorithm}
\usepackage{algpseudocode}
% restore algorithmic keywords so theme doesn't hide them
\algrenewcommand\algorithmicprocedure{\textbf{procedure}}
\algrenewcommand\algorithmicend{\textbf{end procedure}}
\algrenewcommand\algorithmicreturn{\textbf{return}}
\algrenewcommand\textproc[1]{\textnormal{#1}}

% for better lists and spacing control inside algorithms
\usepackage{enumitem}
\setlist{nosep} % remove extra spacing in lists
\algrenewcommand\algorithmicindent{1.0em}

\usepackage{caption}
\captionsetup{labelformat=empty}


%----------------------------------------------------------------------------------------
%    TITLE PAGE
%--------------------------------------------------------------------------------------
\setbeamerfont{title}{size=\large}

\title{NetBERT:Transformer-Driven Robust Intrusion Detection Across Multi-Domain Network Traffic}
\author{\small Project Domain: CyberSecurity\\
Project Guide: Dr.Jyothish K John\\
Team Members: \\
Rhithika M Pradeep -- FIT22CS156 \\
Sreelakshmi P -- FIT22CS180 \\
Ron Binoy Mechery -- FIT22CS159 \\
Sreerag K -- FIT22CS181}
\date{\small May 5, 2026}% Date, can be changed to a custom date

%----------------------------------------------------------------------------------------
%    PRESENTATION SLIDES
%----------------------------------------------------------------------------------------
\usepackage{enumitem}
\setlist[itemize]{label=\textbullet}

\setbeamertemplate{itemize items}[circle]
\begin{document}

\begin{frame}
    % Print the title page as the first slide
    \titlepage
\end{frame}

\section{Abstract}

\begin{frame}{Abstract}
\begin{itemize}

\item Modern networks generate heterogeneous traffic, and encrypted protocols like DoH and DNS tunneling hide malicious activities.

\item Traditional IDS fail to capture contextual relationships and struggle with class imbalance.

\item Proposed \textbf{NetBERT}: a transformer-driven, explainable intrusion detection framework for multi-domain traffic.

\item Uses \textbf{BERT} for contextual feature learning and \textbf{SMOTE–Tomek} for imbalance handling.

\item A lightweight \textbf{MLP classifier} performs intrusion detection.

\item \textbf{LIME-based explainability} and \textbf{Wireshark live packet capture} enable transparent and real-time detection.

\end{itemize}
\end{frame}

\section{Introduction}
\begin{frame}{Introduction}
\begin{itemize}
\item Modern cyberattacks increasingly exploit encrypted DNS channels such as DNS over HTTPS (DoH) and DNS tunneling to bypass traditional Intrusion Detection Systems (IDS).
\item Encryption conceals malicious command-and-control traffic, making payload inspection ineffective for legacy IDS solutions.
\item Transformer-based architectures like BERT extract contextual and semantic relationships from network traffic data, improving detection capability.
\item SMOTomek (SMOTE + Tomek Links) is integrated to address severe class imbalance in intrusion datasets.
\item The SMOTomek–BERT–MLP architecture is extended to encrypted DNS traffic analysis.
\item Explainable AI (XAI) techniques such as LIME are incorporated to ensure transparency and interpretability.
\item Goal: Accurate and explainable detection of DoH, DNS tunneling, and hybrid encrypted attacks, along with non-encrypted attacks.
\end{itemize}
\end{frame}

\section{Problem Statement}
\begin{frame}{Problem Statement}
\begin{itemize}
\item Existing SMOTE–BERT–MLP framework achieves high accuracy but lacks interpretability.
\item They have not been adapted to detect modern encrypted DNS threats,which are harder to analyze.
\item There is a need for an Explainable Deep learning IDS capable of identifying DOH,DNS tunneling, and hybrid attacks while maintaining high accuracy and interpretability.
\end{itemize}
\end{frame}


\section{Project Objectives}
\begin{frame}{Project Objectives}
\begin{itemize}
\item Extend the SMOTETomek–BERT–MLP architecture to the encrypted DNS traffic domain.
\item Integrate Explainable AI methods such as LIME for local interpretability.
\item Improve minority attack classification using enhanced oversampling techniques.
\item Evaluate system performance on datasets such as CIRA-CIC-DoHBrw-2020, UNSW-NB15, CICIDS 2017 and NSL-KDD.
\item Compare enhanced model results with the baseline SMOTE–BERT–MLP framework.
\item Develop explainability dashboards to assist network administrators in interpreting detection decisions.
\item Capture live network and analyze them in-order to classify them as benign or malicious.
\end{itemize}
\end{frame}


\section{Social Relevance}
\begin{frame}{Social Relevance}
\begin{itemize}
\item Strengthens cybersecurity infrastructure against encrypted DNS-based threats.
\item Enhances resilience of smart cities, IoT systems, and digital services.
\item Protects sensitive information from covert data exfiltration channels.
\item Promotes trustworthy AI by integrating interpretability into deep learning-based IDS.
\item Supports global efforts toward secure and reliable digital ecosystems.
\end{itemize}
\end{frame}


\section{SDGs Addressed}
\begin{frame}{SDGs Addressed}
\begin{itemize}

\item \textbf{SDG 9 – Industry, Innovation \& Infrastructure}
\begin{itemize}
    \item Promotes AI-driven innovation in cybersecurity infrastructure.
    \item Strengthens digital networks through intelligent IDS solutions.
\end{itemize}

\item \textbf{SDG 11 – Sustainable Cities \& Communities}
\begin{itemize}
    \item Enhances cybersecurity resilience in smart city and IoT environments.
    \item Ensures secure and privacy-preserving DNS communication.
\end{itemize}

\item \textbf{SDG 16 – Peace, Justice \& Strong Institutions}
\begin{itemize}
    \item Contributes to transparent and explainable AI for accountable cyber defense.
    \item Supports the development of secure and trustworthy online ecosystems.
\end{itemize}

\end{itemize}
\end{frame}


\section{Overview of Literature Review}
\begin{frame}{Overview of Literature Review (1/3)}
\begin{table}[H]
\centering
\setlength{\tabcolsep}{3pt}
\renewcommand{\arraystretch}{1.15}
\caption{Literature Review of Transformer-Based Intrusion Detection Systems}
\resizebox{\textwidth}{!}{
\begin{tabular}{|p{2cm}|p{2.6cm}|p{3.3cm}|p{3cm}|p{4.2cm}|p{3.2cm}|p{3.2cm}|}
\hline
\textbf{Year \& Publisher} &
\textbf{Authors} &
\textbf{Paper Title} &
\textbf{Objectives} &
\textbf{Methodology} &
\textbf{Conclusion / Results} &
\textbf{Limitations} \\ \hline

2024, IEEE Access &
Ali et al. &
BERT Transformer + MLP for Imbalanced IDS &
Improve detection accuracy on imbalanced network traffic &
BERT embeddings + SMOTE + lightweight MLP classifier &
>98\% accuracy on UNSW-NB15, CIC-IDS2017 &
Black-box model; no encrypted DNS handling \\ \hline

2023, ACM DL &
Tripathi et al. &
Flow-based IDS using BERT-MLM &
Model temporal flow relationships &
BERT with Masked Language Modeling treating flows as sentences &
F1 = 0.994 (in-domain), 0.904 (cross-domain) &
High computational cost; no imbalance handling \\ \hline

2022, IEEE Explore &
Zhang \& Bao &
DNN with Connection Efficiency Optimization &
Improve detection efficiency and reduce CPU load &
DNN + Immune Genetic Algorithm + hardware optimization &
High accuracy; reduced processing overhead &
Limited explainability; hardware dependency \\ \hline

2024, ScienceDirect &
Sayyad et al. &
IDS-INT: Transformer Transfer Learning &
Handle class imbalance with transfer learning &
BERT-large + SMOTE + CNN-LSTM hybrid classifier &
99.21\% accuracy; strong generalization &
High training complexity; no XAI integration \\ \hline

\end{tabular}}
\end{table}
\end{frame}

\begin{frame}{Overview of Literature Review (2/3)}
\begin{table}[H]
\centering
\setlength{\tabcolsep}{3pt}
\renewcommand{\arraystretch}{1.15}
\caption{Literature Review of Deep Learning and DNS-Based IDS Methods}
\resizebox{\textwidth}{!}{
\begin{tabular}{|p{2cm}|p{2.6cm}|p{3.3cm}|p{3cm}|p{4.2cm}|p{3.2cm}|p{3.2cm}|}
\hline
\textbf{Year \& Publisher} &
\textbf{Authors} &
\textbf{Paper Title} &
\textbf{Objectives} &
\textbf{Methodology} &
\textbf{Conclusion / Results} &
\textbf{Limitations} \\ \hline

2024, Elsevier &
Anand &
FlowTransformer Framework &
Evaluate shallow vs deep transformers for IDS &
Modular Transformer (BERT, GPT-2); grid search experiments &
F1 up to 98\%; lightweight models feasible &
High computational demand for deep layers \\ \hline

2022, IEEE &
Yin et al. &
Deep Autoencoder-based IDS &
Improve minority attack representation &
Stacked Non-Symmetric Deep Autoencoder (sNDAE) &
>98\% accuracy; low false alarms &
Not suitable for encrypted traffic; no XAI \\ \hline

2024, Springer &
Fernandez et al. &
MLSTL-WSN with SMOTETomek &
Improve imbalance handling in WSN &
Contrastive Learning + SMOTETomek balancing &
Strong generalization across datasets &
Feature distortion risk in synthetic samples \\ \hline

2024, IEEE TIFS &
Gao et al. &
GraphTunnel DNS Detection &
Detect DNS tunneling via graph modeling &
GraphSAGE + graph-to-image + CNN classifier &
99.78\% F1; robust to wildcard tunnels &
Complex graph construction; high memory usage \\ \hline

\end{tabular}}
\end{table}
\end{frame}

\begin{frame}{Overview of Literature Review (3/3)}
\begin{table}[H]
\centering
\setlength{\tabcolsep}{3pt}
\renewcommand{\arraystretch}{1.15}
\caption{Literature Review of Encrypted DNS and Explainable IDS}
\resizebox{\textwidth}{!}{
\begin{tabular}{|p{2cm}|p{2.6cm}|p{3.3cm}|p{3cm}|p{4.2cm}|p{3.2cm}|p{3.2cm}|}
\hline
\textbf{Year \& Publisher} &
\textbf{Authors} &
\textbf{Paper Title} &
\textbf{Objectives} &
\textbf{Methodology} &
\textbf{Conclusion / Results} &
\textbf{Limitations} \\ \hline

2024, IEEE USBEREIT &
Bykov \& Chernyshov &
DNS Tunnel Detection using ML &
Fast and efficient DNS tunnel detection &
LightGBM with entropy \& statistical DNS features &
Macro F1 up to 1.00; low latency &
Limited semantic analysis; no explainability \\ \hline

2025, SpringerOpen &
Mustafa \& Saeed &
BERT + LIME/SHAP for Text Classification &
Improve interpretability of deep models &
BERT embeddings + MLP + LIME/SHAP &
97–99\% accuracy with transparency &
General NLP focus; not IDS-specific \\ \hline

2024, Elsevier &
Alorainy &
Encrypted DNS Detection using BERT &
Detect DoH/DoT encrypted tunnels &
BERT + Graph + SMOTE + MLP fusion &
99.47\% accuracy; low false positives &
Limited cross-domain evaluation \\ \hline

2023, Elsevier &
Ness &
Hybrid CNN–BiLSTM DNS Detection &
Improve adversarial DNS detection &
CNN + BiLSTM + Softmax classifier &
99.12\% accuracy; adversarial robustness &
No XAI integration; limited real-time adaptability \\ \hline

\end{tabular}}
\end{table}
\end{frame}

\section{System Architecture}
\begin{frame}{System Architecture}
\begin{figure}
    \centering
    \includegraphics[width=0.6\linewidth]{system arch.jpeg}
    \label{fig:placeholder}
\end{figure}
\end{frame}

\section*{Project Phases and Results}

% ==========================
% PHASE 1
% ==========================

\begin{frame}{Phases of Project (1/8)}
\textbf{Phase 1: Dataset Collection and Understanding} \\[6pt]

The project utilized benchmark intrusion detection datasets:

\begin{itemize}
    \item CIC-DoHBrw-2020
    \item UNSW-NB15
    \item CICIDS2017
    \item NSL-KDD
\end{itemize}

Each dataset contains labeled benign and malicious network traffic records with multiple attack categories.

Extensive exploratory data analysis (EDA) was performed to understand feature distributions and class imbalance.

\vspace{0.2cm}
\textbf{Outcome:}
Well-understood intrusion datasets prepared for preprocessing.
\end{frame}


% ==========================
% PHASE 2
% ==========================

\begin{frame}{Phases of Project (2/8)}
\textbf{Phase 2: Data Preprocessing and Feature Engineering} \\[6pt]

The following preprocessing steps were applied:

\begin{itemize}
    \item Removal of irrelevant and redundant columns
    \item Handling missing and infinite values
    \item Label encoding of categorical classes
    \item Stratified train-test split
\end{itemize}

To address severe class imbalance:

\begin{itemize}
    \item SMOTETomek technique was applied on training data
\end{itemize}

\vspace{0.2cm}
\textbf{Outcome:}
Balanced and cleaned dataset ready for model training.
\end{frame}


% ==========================
% PHASE 3
% ==========================

\begin{frame}{Phases of Project (3/8)}
\textbf{Phase 3: Feature Transformation using BERT} \\[6pt]

To leverage transformer-based contextual representations:

\begin{itemize}
    \item Tabular network features were converted into structured textual format
    \item Pretrained BERT tokenizer was applied
    \item CLS token embeddings (768-dimensional vectors) were extracted
\end{itemize}

This enabled semantic representation of statistical network traffic features.

\vspace{0.2cm}
\textbf{Outcome:}
High-dimensional contextual embeddings for intrusion detection.
\end{frame}


% ==========================
% PHASE 4
% ==========================

\begin{frame}{Phases of Project (4/8)}
\textbf{Phase 4: Hybrid Model Architecture (BERT + MLP)} \\[6pt]

A hybrid deep learning framework was designed:

\begin{itemize}
    \item Input: 768-dimensional BERT embeddings
    \item Class imbalance handled using SMOTETomek
    \item Multi-Layer Perceptron (MLP) classifier
\end{itemize}

The MLP learns complex nonlinear decision boundaries over contextual embeddings.

\vspace{0.2cm}
\textbf{Outcome:}
Optimized classifier for benign vs malicious traffic prediction.
\end{frame}


% ==========================
% PHASE 5
% ==========================

\begin{frame}{Phases of Project (5/8)}
\textbf{Phase 5: Model Training and Optimization} \\[6pt]

Model training was performed with:

\begin{itemize}
    \item Stratified training-validation split
    \item Adam-based optimization
    \item Early stopping strategy
    \item GPU acceleration (when available)
\end{itemize}

Training loss curves were monitored to ensure convergence and prevent overfitting.

\vspace{0.2cm}
\textbf{Outcome:}
Best-performing trained model checkpoint saved for evaluation.
\end{frame}

% ==========================
% PHASE 6
% ==========================

\begin{frame}{Phases of Project (6/8)}
\textbf{Phase 6: Model Evaluation} \\[6pt]

Performance was evaluated using:

\begin{itemize}
    \item Accuracy
    \item Precision, Recall, F1-Score
    \item Confusion Matrix
    \item ROC-AUC analysis
\end{itemize}

Comparative evaluation was also performed between:

\begin{itemize}
    \item Original tabular features
    \item BERT-transformed embeddings
\end{itemize}

\vspace{0.2cm}
\textbf{Outcome:}
Validated improvement using transformer-based feature learning.
\end{frame}


% ==========================
% PHASE 7
% ==========================

\begin{frame}{Phases of Project (7/8)}
\textbf{Phase 7: Explainable AI using LIME} \\[6pt]

To enhance interpretability:

\begin{itemize}
    \item LIME (Local Interpretable Model-Agnostic Explanations) was applied
    \item Generated local explanations for selected intrusion samples
    \item Identified key contributing features for predictions
\end{itemize}

This provided transparency into model decisions and improved trustworthiness.

\vspace{0.2cm}
\textbf{Outcome:}
Interpretable intrusion detection system with feature-level explanations.
\end{frame}

% ==========================
% PHASE 8
% ==========================

\begin{frame}{Phases of Project (8/8)}
\textbf{Phase 8: Real-Time Intrusion Detection System} \\[6pt]

The trained model was integrated into a live detection framework:

\begin{itemize}
    \item Real-time packet capture using Tshark
    \item Flow-based feature computation
    \item Live BERT embedding generation
    \item Real-time MLP classification
\end{itemize}

Predictions include:

\begin{itemize}
    \item Traffic label (Benign / Malicious)
    \item Confidence score
    \item Automated explanation sentence
\end{itemize}

\vspace{0.2cm}
\textbf{Outcome:}
Fully functional real-time AI-powered intrusion detection system.
\end{frame}

%------------------------ Phase 2 ------------------------
\section{Algorithm}

\begin{frame}{Algorithm (1/3)}
\begin{algorithm}[H]
\small
\begin{algorithmic}
\Procedure{Hybrid\_IDS}{Dataset}

  \State \textbf{Step 1: Dataset Acquisition}
  \State \hspace{1.2em}(a) Load benchmark intrusion datasets
  \State \hspace{1.2em}(b) Identify feature columns and target label
  \State \hspace{1.2em}(c) Perform exploratory data analysis

  \State \textbf{Step 2: Data Preprocessing}
  \State \hspace{1.2em}(a) Remove irrelevant or redundant features
  \State \hspace{1.2em}(b) Handle missing and infinite values
  \State \hspace{1.2em}(c) Encode target labels
  \State \hspace{1.2em}(d) Perform stratified train-test split

\end{algorithmic}
\end{algorithm}
\end{frame}


\begin{frame}{Algorithm (2/3)}
\begin{algorithm}[H]
\small
\begin{algorithmic}

  \State \textbf{Step 3: Feature Transformation using BERT}
  \State \hspace{1.2em}(a) Convert tabular features into structured textual format
  \State \hspace{1.2em}(b) Tokenize text using pretrained BERT tokenizer
  \State \hspace{1.2em}(c) Extract CLS token embeddings (768-dimension)

  \State \textbf{Step 4: Class Imbalance Handling}
  \State \hspace{1.2em}(a) Apply SMOTETomek on training embeddings
  \State \hspace{1.2em}(b) Generate balanced training data

  \State \textbf{Step 5: Model Training}
  \State \hspace{1.2em}(a) Initialize MLP classifier (Adam solver)
  \State \hspace{1.2em}(b) Train on balanced BERT embeddings
  \State \hspace{1.2em}(c) Monitor convergence and loss
  \State \hspace{1.2em}(d) Save trained model
\end{algorithmic}
\end{algorithm}
\end{frame}


\begin{frame}{Algorithm (3/3)}
\begin{algorithm}[H]
\small
\begin{algorithmic}

  \State \textbf{Step 6: Model Evaluation}
  \State \hspace{1.2em}(a) Predict on unseen test data
  \State \hspace{1.2em}(b) Compute Accuracy, Precision, Recall, F1-score
  \State \hspace{1.2em}(c) Generate Confusion Matrix and ROC-AUC

  \State \textbf{Step 7: Explainable AI Integration}
  \State \hspace{1.2em}(a) Apply LIME for local explanation
  \State \hspace{1.2em}(b) Identify influential features per prediction
  \State \hspace{1.2em}(c) Generate human-readable explanation

  \State \textbf{Step 8: Real-Time Deployment}
  \State \hspace{1.2em}(a) Capture live network packets using Tshark
  \State \hspace{1.2em}(b) Compute flow-based statistical features
  \State \hspace{1.2em}(c) Convert features to BERT embedding
  \State \hspace{1.2em}(d) Predict traffic label with confidence score
  \State \hspace{1.2em}(e) Display classification and explanation

\EndProcedure
\end{algorithmic}
\end{algorithm}
\end{frame}

% -------------------------
% End of Algorithms block
% -------------------------
\section{Data Flow Diagrams}
\begin{frame}{Data flow diagrams}
\begin{figure}
    \centering
    \includegraphics[width=0.75\linewidth]{dfd.png}
    \caption{Network Analysis Pipeline}
    \label{fig:placeholder}

\end{figure}
\end{frame}


%---------------------------------- Slide 1 ----------------------------------
\section{Implementation}

\begin{frame}{Implementation -- Tech Stack}
\small
\begin{itemize}

\item \textbf{Language \& Framework:}
Python, PyTorch, HuggingFace Transformers 

\item \textbf{Model Architecture:}
BERT Embeddings (CLS Token), Multi-Layer Perceptron (512-256 Hidden Layers), Batch Normalization, Dropout, Class-Weighted CrossEntropy Loss

\item \textbf{Data \& Preprocessing:}
UNSW-NB15, CIC-DoHBrw-2020, CICIDS2017 and NSL-KDD datasets, Wireshark Packet Capture, Label Encoding, StandardScaler, SMOTETomek (Imbalanced-Learn), Train-Test Split (Stratified)

\item \textbf{Training Optimization:}
AdamW Optimizer, Early Stopping, GPU Acceleration (CUDA), Loss Monitoring, Model Checkpointing

\item \textbf{Evaluation \& Explainability:}
Accuracy, Precision, Recall, F1-Score, Confusion Matrix (Seaborn), LIME (Local Interpretable Model-Agnostic Explanations)

\item \textbf{Deployment \& Inference:}
Real-time User Input Module, BERT-based Feature Embedding Pipeline, Attack  Prediction with Confidence Score



\end{itemize}
\end{frame}



\section{Result}

\begin{frame}{Results (1/10)}
\begin{table}[h]
\centering
\caption{Performance Comparison Across Datasets}
\begin{tabular}{lcccc}
\hline
\textbf{Dataset} & \textbf{Accuracy} & \textbf{Precision} & \textbf{Recall} & \textbf{F1-score} \\
\hline
CIC-DoHBrw-2020 & 0.9827 & 0.98   & 0.98   & 0.98   \\
NSL-KDD         & 0.9893 & 0.9895 & 0.9893 & 0.9894 \\
UNSW-NB15       & 0.9280 & 0.9294 & 0.9280 & 0.9279 \\
CICIDS-2017     & 0.9919 & 0.8663 & 0.8747 & 0.8549 \\
\hline
\end{tabular}
\end{table}
\end{frame}

\begin{frame}{Results (2/10)}
\begin{figure}

    \centering
    \includegraphics[width=0.5\linewidth]{22.jpeg}
    \caption{Confusion matrix(CIC-DoHBrw-2020)}


\end{figure}
\end{frame}

\begin{frame}{Results (3/10)}
\begin{figure}

    \centering
    \includegraphics[width=\linewidth]{ress.png}
    \caption{LIME Explainability(CIC-DoHBrw-2020)}


\end{figure}
\end{frame}

\begin{frame}{Results(4/10)}

\begin{figure}
    \centering
    \includegraphics[width=0.6\linewidth]{NSL-KDD confu.png}
    \caption{Confusion matrix(NSL-KDD Dataset)}
    \label{fig:placeholder}
\end{figure}
\end{frame}

\begin{frame}{Results(5/10)}
\begin{figure}
    \centering
    \includegraphics[width=0.6\linewidth]{NSl-KDD precision.png}
    \caption{Precision-Recall Curves(NSL-KDD Dataset)}
    \label{fig:placeholder}
\end{figure}
\end{frame}

\begin{frame}{Results(6/10)}
\begin{figure}
    \centering
    \includegraphics[width=0.6\linewidth]{NSL-KDD prec-roc.png}
    \caption{Training Loss Curve and ROC Curve(NSL-KDD Dataset)}
    \label{fig:placeholder}
\end{figure}
\end{frame}

\begin{frame}{Results(7/10)}
\begin{figure}
    \centering
    \includegraphics[width=0.5\linewidth]{Confusion NB15.jpg}
    \caption{Confusion Matrix(UNSW-NB15 Dataset)}
    \label{fig:placeholder}
\end{figure}
\end{frame}


\begin{frame}{Results(8/10)}
\begin{figure}
    \centering
    \includegraphics[width=0.6\linewidth]{NB-15 1 new.png}
    \caption{Local Explaination for Class attack and Training loss(UNSW-NB15 Dataset)}
    \label{fig:placeholder}
\end{figure}
\end{frame}


\begin{frame}{Results(9/10)}
\begin{figure}
    \centering
    \includegraphics[width=0.5\linewidth]{CICIDS2017-confusion.jpg}
    \caption{Confusion matrix(CICIDS-2017)}

    \label{fig:placeholder}
\end{figure}
\end{frame}

\begin{frame}{Results(10/10)}
\begin{figure}
    \centering
    \includegraphics[width=0.6\linewidth]{CICIDS2017 roc and prec.png}
    \caption{Precision and ROC Curve(CICIDS-2017)}

    \label{fig:placeholder}
\end{figure}
\end{frame}

\begin{frame}{Testing (1/2)}
\begin{figure}

    \centering
    \includegraphics[width=0.6\linewidth]{Testin 1.png}
    \caption{F1 Score by Attack and Error Rate per Attack}


\end{figure}
\end{frame}

\begin{frame}{Testing (2/2)}
\begin{figure}

    \centering
    \includegraphics[width=0.5\linewidth]{model_comparison_radar_f1_score_NSL-KDD.png}
    \caption{Radar Chart}


\end{figure}
\end{frame}

\section{Comparison with Existing Works}

% -------------------- FRAME 1 --------------------
\begin{frame}{Comparison with Existing Works}
\centering
\scriptsize
\begin{tabular}{|p{3cm}|p{2cm}|p{2cm}|}
\hline
\textbf{Paper} & \textbf{Method} & \textbf{Result} \\ \hline

FlowTransformer &
Transformer (BERT, GPT-2) &
F1 up to 98\% \\ \hline

Deep Learning IDS &
sNDAE &
$>$98\% accuracy \\ \hline

GraphTunnel &
GraphSAGE + CNN &
99.78\% F1 \\ \hline

DNS ML Detection &
LightGBM &
F1 up to 1.00 \\ \hline

NLP + Explainability &
BERT + MLP + LIME &
97–99\% accuracy \\ \hline

Echoes From the Void &
BERT + Graph + SMOTE &
99.47\% accuracy \\ \hline

Hybrid DNS Defense &
CNN + BiLSTM &
99.12\% accuracy \\ \hline

BERT + MLP IDS &
SMOTE + BERT + MLP &
99.3\% accuracy \\ \hline

\textbf{NetBERT (Proposed)} &
SMOTETomek + BERT + MLP + LIME &
\textbf{Improved performance across datasets} \\ \hline


\end{tabular}

\end{frame}

\section{Work Status}

\begin{frame}{Work Status}
\centering
\small
\begin{tabular}{|c|p{6cm}|c|}
\hline
\textbf{Timeline} & \textbf{Activities} & \textbf{Status} \\
\hline

Aug -- Sep & Literature survey on encrypted DNS intrusion detection and transformer-based IDS models & $\checkmark$ \\
\hline

Sep -- Oct & Dataset collection (multi-domain traffic datasets) and preprocessing & $\checkmark$ \\
\hline

Oct -- Nov & BERT model training for contextual feature embedding generation & $\checkmark$ \\
\hline

Nov -- Dec & Implementation of SMOTE–Tomek resampling and MLP classifier training & $\checkmark$ \\
\hline

Dec -- Jan & Integration of LIME-based explainability for model interpretation & $\checkmark$ \\
\hline

Jan -- Feb & Model evaluation on multiple datasets and performance analysis & $\checkmark$ \\
\hline

Feb -- Mar & Integration of real-time packet capture using tshark and pipeline unification & $\checkmark$ \\
\hline

\end{tabular}

\end{frame}


\section{Conclusion}

\begin{frame}{Conclusion}
\begin{itemize}

    \item The system combines \textbf{BERT-based contextual embeddings} with a classical Multi-Layer Perceptron (MLP) classifier to enhance intrusion detection performance.

    \item BERT embeddings capture complex behavioral patterns in both encrypted and unencrypted traffic, improving feature expressiveness beyond traditional numerical representations.

    \item The application of \textbf{SMOTETomek} addresses class imbalance, significantly improving detection accuracy for minority attack categories.

    \item A lightweight \textbf{MLP classifier} ensures efficient multi-class classification while maintaining high predictive performance.

    \item \textbf{LIME-based explainability} enhances transparency by providing interpretable insights into model decisions, increasing trust in predictions.

    \item The integration of a \textbf{real-time flow-based detection module using Tshark} validates the practical deployability of the framework in live network environments.

    \item Overall, the framework presents a \textbf{scalable, interpretable, and hybrid AI-driven intrusion detection solution} suitable for contemporary network infrastructures.

\end{itemize}

\end{frame}

\section{Publication}
\begin{frame}{Publication}

\begin{thebibliography}{00}

\bibitem{}
Dr. Jyothish K John, Rhihtika M Pradeep, Sreelakshmi P, Ron Binoy Mechery, Sreerag K, 
``NetBERT:Transformer-Driven Robust Intrusion Detection Across Multi-Domain
Network Traffic,'' 
International Conference on Recent Trends in Artificial Intelligence, 2025.

\bibitem{}
Dr. Jyothish K John, Rhihtika M Pradeep, Sreelakshmi P, Ron Binoy Mechery, Sreerag K, 
``NetBERT:Transformer-Driven Robust Intrusion Detection Across Multi-Domain
Network Traffic,'' 
2nd International Conference on Innovations in Engineering and Technology, 2026. (Under review)

\end{thebibliography}

\end{frame}

\section{References}
\begin{frame}{References (1/3)}
\small
\begin{thebibliography}{99}

\bibitem{b1}A. Anand, “FlowTransformer: A Transformer Framework for Flow-Based Network Intrusion Detection Systems,” Elsevier, 2024.
\bibitem{b2} Y. Yin et al., “Deep Autoencoder-Based Intrusion Detection System,” IEEE, 2022.

\bibitem{b3} G. Gao, et al., "GraphTunnel: Robust DNS Tunnel Detection Based on DNS Recursive Resolution Graph," \textit{IEEE Transactions on Information Forensics and Security}, vol. 19, pp. 7705-7719, 2024.

\bibitem{b4} N. Bykov and Y. Chernyshov, "Detecting DNS Tunnels Using Machine Learning," in \textit{Proc. 2024 IEEE Ural-Siberian Conf. Biomedical Eng., Radioelectron. Inf. Technol. (USBEREIT)}, pp. 92-94, 2024.

\bibitem{b5} S. Mustafa and M. H. Saeed, "Empowering text classification with NLP and explainable AI for enhanced interpretability," \textit{Journal of Electrical Systems and Information Technology}, vol. 12, no. 81, 2025.

\end{thebibliography}
\end{frame}

\begin{frame}{References (2/3)}
\small
\begin{thebibliography}{99}
\setcounter{enumiv}{5} 

\bibitem{b6} W. S. Alorainy, "Echoes From the Void: Detecting DNS Tunneling With Blackhole Features in Encrypted Scenarios With High Accuracy," \textit{IEEE Access}, vol. 13, pp. 138551-138567, 2025.

\bibitem{b7} S. Ness, "Securing Networks Against Adversarial Domain Name System Tunneling Attacks Using Hybrid Neural Networks," \textit{IEEE Access}, vol. 13, pp. 46697-46709, 2025.


\bibitem{b8} Zeeshan Ali et al. ``Empowering Network Security: BERT Transformer Learning Approach and MLP for Intrusion Detection in Imbalanced Network Traffic'' \emph{IEEE Access}, October 2024.
 A. Ness, “Hybrid CNN-BiLSTM DNS Detection Framework,” Elsevier, 2023.

\bibitem{b9}
Canadian Institute for Cybersecurity (CIC),
"CICIDS2017 Dataset," 2017. [Online]. Available: https://www.unb.ca/cic/datasets/ids-2017.html. [Accessed: Mar. 23, 2026].
\bibitem{b10}
UNSW-NB15 Dataset Paper, "A Comprehensive Dataset for Evaluating Modern Network Intrusion Detection Systems," 2022.

\end{thebibliography}
\end{frame}

\begin{frame}{References (3/3)}
\small
\begin{thebibliography}{99}
\setcounter{enumiv}{10} 

\bibitem{b11}
UCI Machine Learning Repository,
"KDD Cup 1999 Data," 1999. [Online]. Available: http://kdd.ics.uci.edu/databases/kddcup99/kddcup99.html. [Accessed: Mar. 23, 2026].
\bibitem{b12}
Dhoogla,
"CIC-DoHBrw-2020 Dataset," Kaggle, 2021. [Online]. Available: https://www.kaggle.com/datasets/dhoogla/cicdohbrw2020. [Accessed: Mar. 23, 2026].

\bibitem{b13}
N. Tripathi \textit{et al.},
``Flow-based network intrusion detection based on BERT masked language model,''
\emph{Proceedings of the ACM Digital Library}, 2023.

\bibitem{b14}
T. Zhang and S. Bao,
``A novel deep neural network model for computer network intrusion detection,''
\emph{IEEE Access / IEEE Explore}, 2022.

\bibitem{b15}
F. Sayyad \textit{et al.},
``IDS-INT: Intrusion detection system using transformer-based transfer learning,''
\emph{ScienceDirect (Digital Communications and Networks)}, 2024.

\bibitem{b16}
A. Fernández \textit{et al.},
``MLSTL-WSN: Machine learning-based intrusion detection using SMOTETomek in wireless sensor networks,''
\emph{Neural Computing and Applications}, Springer, 2024.

\end{thebibliography}
\end{frame}




\section{Contribution of Each Member}

\begin{frame}{Contribution of Each Member}

\begin{table}[h!]
\centering
\renewcommand{\arraystretch}{1.3}
\small
\begin{tabular}{|p{3cm}|p{9.5cm}|}
\hline
\textbf{Member} & \textbf{Contributions} \\ \hline

\textbf{Sreerag K} &
Implemented the hybrid BERT + SMOTETomek + MLP intrusion detection pipeline on the CIC-DoH dataset, including preprocessing, evaluation, LIME-based explainability, and real-time detection using Tshark. Prepared presentation slides.
\\ \hline

\textbf{Ron Binoy Mechery} &
Implemented the hybrid BERT + SMOTETomek + MLP intrusion detection pipeline on the NSL-KDD dataset, including preprocessing, evaluation, explainability, and real-time deployment. Prepared presentation slides.
\\ \hline

\textbf{Rhithika M Pradeep} &
Implemented the hybrid BERT + SMOTETomek + MLP intrusion detection pipeline on the CICIDS2017 dataset, including preprocessing, evaluation, explainability, and real-time deployment. Prepared presentation slides.
\\ \hline

\textbf{Sreelakshmi P} &
Implemented the hybrid BERT + SMOTETomek + MLP intrusion detection pipeline on the UNSW-NB15 dataset, including preprocessing, evaluation, explainability, and real-time deployment. Prepared presentation slides.
\\ \hline

\end{tabular}
\end{table}

\end{frame}


\end{document}